\expandafter\let\csname Ecal\endcsname\undefined
\expandafter\let\csname Hil\endcsname\undefined
\expandafter\let\csname Sch\endcsname\undefined
\expandafter\let\csname bumppath\endcsname\undefined
\expandafter\let\csname card\endcsname\undefined
\expandafter\let\csname diam\endcsname\undefined
\expandafter\let\csname dist\endcsname\undefined
\expandafter\let\csname esssup\endcsname\undefined
\expandafter\let\csname id\endcsname\undefined
\expandafter\let\csname ip\endcsname\undefined
\expandafter\let\csname sgn\endcsname\undefined
\expandafter\let\csname spn\endcsname\undefined
\expandafter\let\csname supp\endcsname\undefined
\expandafter\let\csname weakto\endcsname\undefined
\expandafter\let\csname wsto\endcsname\undefined
\expandafter\let\csname wt\endcsname\undefined
\newcommand{\ip}[2]{\langle #1, #2 \rangle}
\newcommand{\spn}{\operatorname{span}}
\newcommand{\wt}{\widetilde}
\newcommand{\id}{\operatorname{id}}
\newcommand{\Sch}{\mathcal{S}}
\newcommand{\Ecal}{\mathcal{E}}
\newcommand{\Hil}{\mathcal{H}}
\newcommand{\weakto}{\rightharpoonup}
\newcommand{\wsto}{\overset{*}{\rightharpoonup}}
\newcommand{\supp}{\operatorname{supp}}
\newcommand{\esssup}{\operatorname*{ess\,sup}}
\newcommand{\sgn}{\operatorname{sgn}}
\newcommand{\dist}{\operatorname{dist}}
\newcommand{\diam}{\operatorname{diam}}
\newcommand{\card}{\operatorname{card}}

\newcommand{\bumppath}{plot coordinates {
	(-1,0) (-0.95,0.0001) (-0.9,0.0141) (-0.85,0.0740) (-0.8,0.1691)
	(-0.75,0.2657) (-0.7,0.3825) (-0.65,0.4788) (-0.6,0.5698)
	(-0.55,0.6482) (-0.5,0.7165) (-0.45,0.7756) (-0.4,0.8266)
	(-0.35,0.8700) (-0.3,0.9057) (-0.25,0.9349) (-0.2,0.9594)
	(-0.15,0.9779) (-0.1,0.9900) (-0.05,0.9975) (0,1)
	(0.05,0.9975) (0.1,0.9900) (0.15,0.9779) (0.2,0.9594)
	(0.25,0.9349) (0.3,0.9057) (0.35,0.8700) (0.4,0.8266)
	(0.45,0.7756) (0.5,0.7165) (0.55,0.6482) (0.6,0.5698)
	(0.65,0.4788) (0.7,0.3825) (0.75,0.2657) (0.8,0.1691)
	(0.85,0.0740) (0.9,0.0141) (0.95,0.0001) (1,0)}}
